% Conclusion

We have demonstrated that edge-weighted recursive bisection redistricting transfers successfully to six parliamentary democracies, achieving compactness improvements of 10--22\% over enacted boundaries across FPTP, STV, and MMP electoral systems. For multi-member systems (Ireland, Malta), geographic optimization does not impair proportionality: D'Hondt seat allocation within compact, geographically coherent constituencies achieves Gallagher indices below 2.0 regardless of boundary-drawing method.

Several implications follow. First, algorithmic redistricting is not a solution to a uniquely American problem. The underlying tension between geographic integrity and political manipulation exists in any system where boundaries affect electoral outcomes --- including systems with independent boundary commissions. Second, multi-member PR systems are more robust to redistricting manipulation than single-member systems: the proportional allocation mechanism within each constituency absorbs geographic variation in vote shares. This suggests that multi-member systems may reduce the stakes of redistricting disputes, with implications for US electoral reform proposals.

Third, cross-national redistricting analysis requires vocabulary adaptation beyond algorithm adaptation. The RPLAN open format and \texttt{redist policy} database introduced in this paper provide a foundation for systematic cross-national redistricting research, enabling future work to compare not just algorithmic outcomes but the institutional determinants of redistricting quality.

Future work should extend this analysis to three additional contexts: (1) subnational redistricting in federal systems (German Länder, Canadian provinces, Australian states); (2) indigenous electoral districts (New Zealand Māori electorates, Canadian First Nations constituencies); and (3) the transition from single-member to multi-member systems, using algorithmic redistricting to model the boundary implications of electoral reform proposals. The RPLAN format and open redistricting platform make these extensions tractable for the first time.
